\begin{figure*}[t!]
  \centering
  \includegraphics[width = 1.0\textwidth]
  {./fig/main_method_iccv.pdf}
  \vspace{-2em}
    \caption{\textbf{Overview of Scene Graph Guided 3D Scene Generation.} The Scene Graph Guided 3D Generation structure consists of three main components: the interactive system (red), BEM processing (blue), and diffusion generation (bottom). Through the interactive system, users can construct their own Scene Graphs using either an interactive interface or text interaction. The constructed scene graph is processed by a GNN, which is jointly trained with the diffusion model using auxiliary tasks to enhance control. Each node in the Scene Graph is then positioned by the Allocation Module to form the BEM. This BEM serves as a conditioning input to the 3D Pyramid Discrete Diffusion Model~\cite{pdd}, which generates the final 3D outdoor scene. Note that ``Recon", ``Classification", and ``CANE" denote ``Edge Reconstruction", ``Node Classification", and ``Context-aware Node Embedding", respectively.}
    \label{fig:main_method}
    % \vspace{-1em}
    \vspace{1em}
\end{figure*}

\section{Related Work}
\label{sec:related_work}

\noindent \textbf{3D Generation via Diffusion Models.} 
Diffusion models have expanded their applications from 2D image synthesis to complex 3D data modeling~\cite{pamiDMiV2023}. 
%
Compared to conventional GANs~\cite{GAN} and VAEs~\cite{VAE}, they offer improved performance through a progressive denoising mechanism~\cite{DDPM}, enhancing training stability and the ability to model complex distributions.
%
This makes them particularly effective for 3D data generation. 
%
However, existing research primarily focuses on object-level generation~\cite{xu2025sparp, ren2024xcube, magic123, liu2023zero1to3, yang2019pointflow, xiang2024structured} or indoor environments~\cite{Bokhovkin2024SceneFactorFL, tang2024diffuscene, Yang2024SceneCraftL3, Fang2023CtrlRoomCT}. 
%
Meanwhile, the few works~\cite{lin2023infinicity, xie2024citydreamer,pdd,lee2024semcity,lee2023diffusion} designed for outdoor scene generation prioritize visual fidelity over controllability. 
%
In this work, we aim to develop a framework for controllable 3D outdoor scene generation, emphasizing easy and precise user control.

\noindent \textbf{Scene Graph Application.} 
A scene graph is a structured representation of a scene, encoding objects, their attributes, and the relationships between them~\cite{Johnson_2015_CVPR}. 
%
Unlike dense representations such as point clouds~\cite{qi2017pointnet++,li2025deephierarchicallearningfor3dsemanticsegmentation,li2024enhancing} or meshes~\cite{chang2015shapenet}, a scene graph provides a concise yet comprehensive view of a scene’s structure, making it a powerful conditional signal for generation tasks. 
%
By explicitly modeling inter-object relationships, scene graphs offer a structured way to control scenes, facilitating both human-driven and AI-driven content generation, and this capability has been widely explored especially in 2D tasks~\cite{SonSGpami2023, Liu2024R3CDSG, yang2022dbsgi, Yang_2019_CVPR, XU2019477, vcg2021tmm, johnson2018image, farshad2023scenegenie}. 
%
Extending this concept to 3D environments, Armeni et al.~\cite{armeni_iccv19} introduce the 3D indoor scene graph. This framework integrates semantic, spatial, and geometric information into a hierarchical graph, where nodes represent objects, rooms, and spaces, and edges encode their spatial and semantic relationships. Although this well-defined representation has proven effective for indoor scenes, a comparable formulation for outdoor environments remains largely unexplored.

\noindent \textbf{Controllable 3D Scene Generation.} 
%
To date, controllable 3D outdoor scene generation has received limited attention. One of the closest efforts in this area is Text2LiDAR~\cite{wu2024text2lidar}, which generates LiDAR points from text inputs.
%
While text-based control is an interesting attempt, it lacks explicit spatial structure, making it unsuitable for precise scene composition~\cite{zhang2023adding}.
%
Instead, scene graphs offer a more interpretable and structured approach to controlling scene generation, a concept that has been well established in indoor scene generation~\cite{lin2024instructscene, tang2024diffuscene, zhou2019scenegraphnet, li2019grains, wang2021sceneformer}. 
%
However, adapting them to outdoor scenes is non-trivial, as outdoor scenes are large-scale, unbounded, and contain a diverse set of objects~\cite{xie2025citydreamer4dcompositionalgenerativemodel}. 
%
Unlike indoor scenes, which often use a compositional approach by placing objects within predefined or generated bounding boxes~\cite{zhai2025echoscene, zhai2023commonscenes, dhamo2021graph, wang2019planit, para2021generative}, outdoor scenes feature complex backgrounds and incomplete structures, such as roads and buildings.
%
To address these challenges, we propose a novel pipeline and scene graph representation tailored for controllable 3D outdoor scene generation.






